

Inicia-se com a versão forte da equação de equilíbrio mecânico para o caso de uma viga engastada em uma extremidade:
\begin{figure}
    \centering
    \includegraphics[width=0.5\linewidth]{image.png}
    \caption{Viga engastada}
    \label{Viga}
\end{figure}
Temos então o PVC:

\begin{equation}
   \left\{\begin{aligned}
        -\nabla\cdot\sigma = \textbf{f}, \quad \mathbf{x}\in\Omega,\\
        \sigma = \lambda tr(\epsilon)I + 2\mu\epsilon,\\
        \epsilon = \frac{1}{2}(\nabla \mathbf{u} + \nabla \mathbf{u}^T),\\
        \mathbf{u} = 0, \mathbf{x}\in\Gamma_u,\\
        \mathbf{\sigma n} = 0, \mathbf{x}\in\Gamma_t,
    \end{aligned}
    \right.
\end{equation}

onde $\Gamma_u$,  é a superfície onde são aplicadas as condições de contorno de Dirichlet e  $\Gamma_t$ aquela onde são aplicadas as condições de contorno de Neumann. Para encontrar a forma fraca da equação de equilíbrio, iniciamos multiplicando ambos os lados da equação por uma função vetorial de teste $\mathbf{v}$, com o mesmo formato que o campo de deslocamentos.

\begin{equation}
 -(\nabla\cdot\sigma)^T\mathbf{v} = \textbf{f}^T\mathbf{v}.
\end{equation}
Integra-se ambos os lados,
\begin{}

\begin{equation}
 -\int_{\Omega}(\nabla\cdot\sigma)^T\mathbf{v}d\Omega =\int_{\Omega} \textbf{f}^T\mathbf{v}d\Omega.
\end{equation}
O próximo passo é aplicar a integração por partes do lado esquerdo. Para isso, devemos considerar o seguinte,

\begin{equation}
    \nabla\cdot(\sigma\mathbf{v}) = (\nabla\cdot\sigma)^T\mathbf{v} + \sigma:\nabla\mathbf{v},
\end{equation}
\begin{equation}
- (\nabla\cdot\sigma)^T\mathbf{v} =  \sigma:\nabla\mathbf{v}-\nabla\cdot(\sigma\mathbf{v})   ,
\end{equation}
onde $\sigma:\nabla\mathbf{v}$ indica a operação de contração dupla do tensor de tensões pelo gradiente da função teste. Retornando à integral,

\begin{equation}
 \int_{\Omega}  \sigma:\nabla\mathbf{v} d\Omega- \int_{\Omega}\nabla\cdot(\sigma\mathbf{v}) d\Omega =\int_{\Omega} \textbf{f}^T\mathbf{v}d\Omega.
\end{equation}
Pelo teorema da divergência de Gauss, é possível reduzir a ordem do termo $\int_{\Omega}\nabla\cdot(\sigma\mathbf{v}) d\Omega$,

\begin{equation}
    \int_{\Omega}\nabla\cdot(\sigma\mathbf{v})d\Omega= \int_{\partial\Omega}(\sigma\mathbf{v})^T\mathbf{n}dS
\end{equation}
Considerando que o tensor $\sigma$ é simétrico, o produto $(\sigma\mathbf{v})^T\mathbf{n}$ pode ser desenvolvido da seguinte forma,

\begin{equation}
    (\sigma\mathbf{v})^T\mathbf{n} = \sigma^T\mathbf{v}^T\mathbf{n} =  \mathbf{v}^T\sigma\mathbf{n} = \mathbf{v}^T\mathbf{t}
\end{equation}
onde $\mathbf{t}$ é o vetor de tração, produto da aplicação de uma tensão a uma determinada área. Outro ponto a ser considerado é a contração dupla $\sigma:\nabla\mathbf{v}$, onde $\nabla\mathbf{v}$ pode ser separado em sua parte simétrica e anti-simétrica, ou seja,

\begin{equation}
    \sigma:\nabla\mathbf{v} =  \sigma:sym(\nabla\mathbf{v}) + \sigma:antisym(\nabla\mathbf{v}) = \sigma:\left[\frac{1}{2}(\nabla\mathbf{v} + \nabla\mathbf{v}^T)\right] + \sigma:\left[\frac{1}{2}(\nabla\mathbf{v} - \nabla\mathbf{v}^T)\right],
\end{equation}
Mas, sabe-se da álgebra linear que a contração dupla de um tensor simétrico por um anti-simétrico é nula. Logo, apenas a parte simétrica de $\mathbf{v}$ prevalece.  Volta-se então ao problema original,

\begin{equation}
 \int_{\Omega}  \sigma(\mathbf{u}):\epsilon(\mathbf{v}) d\Omega  =\int_{\Omega} \textbf{f}^T\mathbf{v}d\Omega + \int_{\partial\Omega}\mathbf{t}^T\mathbf{v}dS,
\end{equation}
com $\epsilon(\mathbf{v})$, indicando a parte simétrica do gradiente do campo $\mathbf{v}$.

